\chapter{Background and Foundations}
\label{ch:background}

\section{Large Language Models and Alignment}

Large language models are first trained to model natural-language distributions and are then commonly post-trained to better follow instructions and satisfy behavioral preferences. Reinforcement learning from human feedback (RLHF) is a prominent example: demonstrations and ranked outputs are used to train a preference signal that later guides policy optimization \parencite{ouyang2022instructgpt}. The important point for this thesis is that alignment is mediated by operational signals. Demonstrations, constitutions, human preferences, policies or evaluators define what behavior is rewarded.

Such signals are necessarily incomplete. A model can be generally helpful while still producing advice that is inappropriate for the explicit cultural, linguistic, legal or interpersonal context of a request. Conversely, disagreement with a population majority is not automatically culturally inappropriate. Pluralistic-alignment research therefore argues that systems serving heterogeneous populations may need to represent or accommodate multiple reasonable perspectives rather than collapse all judgments into one global preference \parencite{sorensen2024pluralistic}.

This thesis uses cultural alignment in an operational sense: a response should fit the explicit situation, avoid materially misrepresenting relevant practices or constraints, and preserve legitimate variation when the context does not justify a single categorical rule.

\section{Preference Alignment and Reward Models}

Reward models transform preference data into a reusable evaluation signal. A classifier RM often assigns a scalar reward to a prompt--response pair, while generative evaluators produce scores or judgments through language generation. General RM benchmarks such as RewardBench measure preference discrimination rather than cultural evidence gathering \parencite{lambert2025rewardbench}. Skywork Reward V2, used later as an independent baseline, is a general preference-model family developed through large-scale preference-data curation \parencite{liu2026skyworkrewardv2}.

A scalar reward is useful for ranking but limited diagnostically. It does not reveal which cultural claim, pragmatic convention, institutional constraint or stereotype caused the ranking. It is also a learned preference signal rather than a calibrated probability of cultural correctness. For this reason, the reward model remains strictly external to Vericult: RM scores never enter Vericult retrieval, scoring or ranking.

\section{LLM-as-a-Judge}

A second evaluation paradigm prompts an instruction-tuned LLM directly to compare or rate responses. LLM-as-a-judge evaluation is attractive because it is flexible and can express criteria in natural language. However, strong judges still exhibit systematic effects. Work on MT-Bench and related settings reports position, verbosity and self-enhancement biases, while later analyses show that response ordering can materially change pairwise judgments \parencite{zheng2023llmjudge,wang2024notfairevaluators}.

This thesis therefore treats a direct judge as a comparator rather than as the verifier. The direct judge sees the candidate responses and decides among them. Vericult instead decomposes the response, retrieves candidate-blind evidence where appropriate, and preserves the evidence path. Whether this additional structure helps is an empirical question rather than an assumption.

\section{Cultural Appropriateness}
\label{sec:cultural-appropriateness}

The term \emph{cultural correctness} suggests a fixed, homogeneous cultural truth. This is too strong for many real-world questions. UNESCO describes culture through material, intellectual and emotional features, lifestyles, value systems, traditions and beliefs, while the UN Committee on Economic, Social and Cultural Rights treats cultural life as broad, evolving and connected to identity, language, knowledge, customs, institutions and participation \parencite{unesco2001culturaldiversity,cescr2009gc21}. Surveys of cultural NLP likewise note that research operationalizes selected proxies and facets rather than one universally accepted computational definition \parencite{adilazuarda2024culture,liu2025culturallyaware,pawar2025culturalawareness}.

Accordingly, this thesis uses \emph{cultural appropriateness} for a context-sensitive evaluative construct. Several epistemic distinctions matter:

\paragraph{Descriptive tendencies.} Evidence that a practice is common does not establish a universal rule. Cross-national opinion data can show population differences while preserving substantial within-population variation \parencite{durmus2023globalopinionqa,tao2024culturalbias}.

\paragraph{Social norms and etiquette.} Social expectations vary with relationship, setting, generation and locality. NormAd demonstrates how model judgments change as cultural norms become more explicit \parencite{rao2025normad}.

\paragraph{Pragmatic conventions.} Register, forms of address, indirectness, idioms and implicature affect cultural fit even when grammar and literal factuality are correct.

\paragraph{Values.} Values are plural and contested. Cultural evaluation should not transform majority preferences into moral unanimity.

\paragraph{Law and institutional rules.} Formal obligations differ from informal practice. SafeWorld demonstrates the need to consider both culturally situated norms and geographically varying legal constraints \parencite{yin2024safeworld}.

These distinctions motivate a verifier that matches the strength and kind of its conclusion to the available context and evidence.

\section{Evidence-Grounded Evaluation and Retrieval}

Retrieval-augmented generation established the broader principle that a language model can combine parametric knowledge with explicitly retrieved external information \parencite{lewis2020rag}. Vericult does not use retrieval to generate the user's answer; it uses retrieval to make parts of the evaluation inspectable. This is particularly valuable for claims about law, institutions, current policy, regional practice or the prevalence of cultural conventions.

Retrieval is not a source of ground truth. Search availability is uneven across languages and communities; provider snippets can omit qualifications; and source authority depends on the question. An official source is appropriate for a legal rule, while community material can be informative for lived practice. Vericult therefore uses source classes as provenance descriptors rather than a universal numeric hierarchy.

Retrieval can also contaminate evaluation if benchmark answers, annotations or result files are encountered. Benchmark-contamination research motivates explicit exclusions and artifact recording \parencite{deng2024contamination}. The implementation filters known repository and annotation paths and archives the exact evidence used, while acknowledging that copied or mirrored material cannot be eliminated perfectly.

\section{Evaluation Concepts}

\paragraph{Human reference versus ground truth.} Cultural annotations provide an operational comparison target, not objective universal truth. Agreement among annotators is itself part of the result.

\paragraph{Absolute and relative evaluation.} A candidate can be the best among four poor responses. Candidate-level appropriateness and Best-of-4 preference therefore answer different questions and must be analyzed separately.

\paragraph{Evidence uncertainty.} A completed Vericult run always returns a cultural label, but evidence sufficiency and contextual-fallback usage remain separate diagnostics. A label supported by strong retrieval evidence is not methodologically identical to one produced after unresolved retrieval.

\paragraph{Paired comparison.} Because the same items are evaluated by Vericult, the reward model and the direct judge, system comparisons are paired at prompt level.

\paragraph{Reproducibility.} Fixing temperature is insufficient. Model versions, prompt templates, retrieval dates, evidence, ordering and code revisions can change outcomes. Vericult therefore distinguishes LIVE evidence acquisition from REPLAY evaluation over frozen retrieval artifacts.

These foundations motivate the research gap addressed in the next chapter: existing evaluators provide valuable preference and capability measurements, but a practical cultural verifier should additionally expose what is being judged, what external support was found and where its decision came from.
